\documentclass[troisieme]{fiche}
\metadonnees{2}{Le Kansas gris}{Bloc A — Raconter}

\begin{document}
\entete

\objectifs{%
\textbf{Langue} : le prétérit ; verbes irréguliers, lot 1.\par
\textbf{Texte} : L. Frank Baum, \emph{The Wonderful Wizard of Oz} (1900), ch. \textsc{i}.\par
\textbf{Durée} : 1 h — exercices 20 min, texte et questions 25 min, version 15 min.}

%% ===================================================================
\exercice[12 min]{Le prétérit}

\rappel{\textbf{Rappel.} Le prétérit raconte ce qui est fini. Les verbes réguliers prennent
\emph{-ed} ; les autres ont une forme à apprendre.}

\consigne{Mettez le verbe entre parenthèses au prétérit.}

\begin{anglais}
\begin{enumerate}[label=\arabic*.,itemsep=2.5mm,leftmargin=*]
  \item Dorothy \emph{(live)} \trou{lived} in the middle of the great Kansas prairies.
  \item Their house \emph{(be)} \trou{was} small, with only one room.
  \item Uncle Henry \emph{(work)} \trou{worked} hard from morning till night.
  \item He never \emph{(laugh)} \trou{laughed} and rarely \emph{(speak)} \trou{spoke}.
  \item The sun \emph{(burn)} \trou{burned \emph{ou} burnt} the grass until it turned gray.
  \item Aunt Em \emph{(come)} \trou{came} there as a young, pretty wife.
  \item The wind and the sun \emph{(take)} \trou{took} the red from her cheeks.
  \item Toto \emph{(play)} \trou{played} all day long and Dorothy \emph{(love)} \trou{loved}
    him dearly.
  \item Dorothy \emph{(see)} \trou{saw} nothing but the great gray prairie.
  \item The family \emph{(go)} \trou{went} down into the cyclone cellar.
\end{enumerate}
\end{anglais}

\note[Orthographe du \emph{-ed}]{Consonne + \emph{y} devient \emph{-ied} :
\emph{carry $\rightarrow$ carried}, mais \emph{play $\rightarrow$ played}. Une consonne
finale unique précédée d'une voyelle accentuée double : \emph{stop $\rightarrow$ stopped},
\emph{plan $\rightarrow$ planned}.}

\note[Irréguliers du texte]{\emph{be $\rightarrow$ was / were}, \emph{come $\rightarrow$
came}, \emph{take $\rightarrow$ took}, \emph{see $\rightarrow$ saw}, \emph{go $\rightarrow$
went}, \emph{speak $\rightarrow$ spoke}. \emph{Burn} accepte les deux formes :
\emph{burned} ou \emph{burnt}, comme \emph{learn}, \emph{dream}, \emph{spell}.}

%% ===================================================================
\exercice[8 min]{Verbes irréguliers, lot 1}
\consigne{a) Complétez. À savoir par cœur pour la séance prochaine.}

\begin{center}\small
\begin{tabular}{@{}lll@{}}
\textsf{base} & \textsf{prétérit} & \textsf{participe passé} \\[1.5mm]
\emph{be} & \trou{was / were} & \trou{been} \\[1mm]
\emph{have} & \trou{had} & \trou{had} \\[1mm]
\emph{do} & \trou{did} & \trou{done} \\[1mm]
\emph{go} & \trou{went} & \trou{gone} \\[1mm]
\emph{come} & \trou{came} & \trou{come} \\[1mm]
\emph{see} & \trou{saw} & \trou{seen} \\[1mm]
\emph{say} & \trou{said} & \trou{said} \\[1mm]
\emph{make} & \trou{made} & \trou{made} \\[1mm]
\emph{take} & \trou{took} & \trou{taken} \\[1mm]
\emph{get} & \trou{got} & \trou{got \emph{ou} gotten} \\
\end{tabular}
\end{center}

\consigne{b) Complétez avec un verbe du tableau, au prétérit.}
\begin{anglais}
\begin{enumerate}[label=\arabic*.,itemsep=2.5mm,leftmargin=*]
  \item Aunt Em \trou{came} to the prairie when she was young.
  \item Dorothy \trou{saw} only gray grass and a gray sky.
  \item The little dog \trou{made} her laugh every day.
  \item Uncle Henry \trou{said} very little and \trou{did} his work in silence.
\end{enumerate}
\end{anglais}

\note[Trois remarques]{\emph{Said} et \emph{made} s'écrivent avec un seul \emph{d} ; on
entend \og sed \fg{} et non \og séd \fg. \emph{Come} et \emph{become} ont le même participe
passé que la base. \emph{Gotten} est américain ; en anglais britannique, \emph{got} suffit.}

%% ===================================================================
\newpage
\section{Texte}

\noindent\small\textit{Dorothy est orpheline. Elle vit avec son oncle et sa tante dans une
ferme isolée des grandes plaines du Kansas.}\normalsize

\begin{texte}
When Dorothy stood in the doorway and looked around, she could see nothing but the great gray
\gl[a prairie]{prairie}{une prairie, une grande plaine} on every side. Not a tree nor a house
broke the broad \gl[a sweep]{sweep}{une étendue} of flat country that reached to the edge of
the sky in all directions. The sun had baked the \gl[to plow]{plowed}{labourer} land into a
gray mass, with little \gl[a crack]{cracks}{une fissure} running through it. Even the grass
was not green, for the sun had burned the tops of the long \gl[a blade]{blades}{un brin
(d'herbe)} until they were the same gray color to be seen everywhere. Once the house had been
painted, but the sun \gl[to blister]{blistered}{faire cloquer} the paint and the rains washed
it away, and now the house was as \gl[dull]{dull}{terne} and gray as everything else.

When Aunt Em came there to live she was a young, pretty wife. The sun and wind had changed
her, too. They had taken the \gl[a sparkle]{sparkle}{un éclat} from her eyes and left them a
\gl[sober]{sober}{sobre, éteint} gray; they had taken the red from her cheeks and lips, and
they were gray also. She was thin and \gl[gaunt]{gaunt}{décharné}, and never smiled now. When
Dorothy, who was an orphan, first came to her, Aunt Em had been so startled by the child's
laughter that she would scream and press her hand upon her heart whenever Dorothy's
\gl[merry]{merry}{joyeux} voice reached her ears; and she still looked at the little girl with
wonder that she could find anything to laugh at.

Uncle Henry never laughed. He worked hard from morning till night and did not know what joy
was. He was gray also, from his long beard to his rough boots, and he looked
\gl[stern]{stern}{sévère} and \gl[solemn]{solemn}{grave}, and rarely spoke.

It was Toto that made Dorothy laugh, and saved her from growing as gray as her other
surroundings. Toto was not gray; he was a little black dog, with long silky hair and small
black eyes that \gl[to twinkle]{twinkled}{pétiller} merrily on either side of his funny,
\gl[wee]{wee}{tout petit} nose. Toto played all day long, and Dorothy played with him, and
loved him dearly.
\end{texte}
\source{L. Frank Baum, \emph{The Wonderful Wizard of Oz}, Chicago, Hill, 1900,
ch.~\textsc{i}.}

\partiel[15 min]{Questions}
\consigne{Répondez aux deux premières questions en français, aux deux dernières en anglais.}

\begin{enumerate}[label=\arabic*.,itemsep=2.5mm,leftmargin=*]
  \item Combien de fois le mot \emph{gray} apparaît-il dans le texte, et à propos de quoi ?
    \lignes{2}
    \corr{Sept fois : la prairie, la terre labourée, l'herbe, la maison, les yeux de tante
    Em, ses joues et ses lèvres, oncle Henry tout entier. Tout ce que le soleil a touché est
    devenu gris.}
  \item Qu'est-ce qui, dans ce paysage, n'est pas gris ? Qu'est-ce que cela change ?
    \lignes{2}
    \corr{Toto, le petit chien noir aux yeux pétillants. C'est lui qui fait rire Dorothy et
    qui l'empêche de devenir grise comme le reste — la seule chose vivante du texte.}
  \item \en{What had the sun and the wind done to Aunt Em?}
    \lignes{2}
    \corr{They had taken the sparkle from her eyes and the red from her cheeks and lips, and
    left everything gray. She had become thin and gaunt, and she never smiled.}
  \item \en{Find in the text three verbs in the preterit and give their base form.}
    \lignes{2}
    \corr{Par exemple \emph{stood} (stand), \emph{broke} (break), \emph{came} (come),
    \emph{took} (take), \emph{made} (make), \emph{spoke} (speak) ; et les réguliers
    \emph{looked}, \emph{reached}, \emph{washed}, \emph{played}, \emph{loved}.}
\end{enumerate}

\note[Une forme en avance]{\emph{The sun had baked}, \emph{they had taken}, \emph{the house
had been painted} : ce n'est plus le prétérit mais le \emph{pluperfect}, qui remonte plus
loin dans le passé. On l'étudie en seconde ; ici il suffit de le reconnaître et de le
traduire par un plus-que-parfait.}

%% ===================================================================
\section{Version}
\consigne{Traduisez le passage en français. C'est le début du chapitre, juste avant le
texte.}

\begin{version}
Dorothy lived in the midst of the great Kansas prairies, with Uncle Henry, who was a farmer,
and Aunt Em, who was the farmer's wife. Their house was small, for the
\gl[lumber]{lumber}{le bois de construction} to build it had to be carried by
\gl[a wagon]{wagon}{un chariot} many miles. There were four walls, a floor and a roof, which
made one room; and this room contained a rusty looking
\gl[a cookstove]{cookstove}{une cuisinière}, a \gl[a cupboard]{cupboard}{un placard} for the
dishes, a table, three or four chairs, and the beds.
\end{version}
\source{\emph{Ibid.}}

\lignes{10}

\corr{\textbf{Proposition de traduction.} Dorothy vivait au milieu des grandes plaines du
Kansas, avec l'oncle Henry, qui était fermier, et la tante Em, qui était la femme du fermier.
Leur maison était petite, car il avait fallu charrier sur des kilomètres le bois nécessaire à
sa construction. Il y avait quatre murs, un plancher et un toit, qui formaient une seule
pièce ; et cette pièce contenait une cuisinière d'aspect rouillé, un placard pour la
vaisselle, une table, trois ou quatre chaises, et les lits.}

\note[Choix de traduction 1]{\emph{had to be carried} : \emph{have to} exprime l'obligation
(\og il fallait \fg) et le passif se rend par une tournure active avec \og il \fg{}
impersonnel : \og il avait fallu charrier \fg. Traduire mot à mot donnerait \og devait être
porté \fg, correct mais lourd.}

\note[Choix de traduction 2]{\emph{a rusty looking cookstove} : \emph{-looking} forme des
adjectifs composés, \og d'aspect rouillé \fg, \og qui a l'air rouillé \fg. De même
\emph{good-looking}, \emph{strange-looking}.}

\note[Choix de traduction 3]{\emph{many miles} : un mile fait 1,6 km. On peut garder
\og des miles \fg, mais \og sur des kilomètres \fg{} passe mieux en français, l'idée étant la
distance et non la mesure exacte.}

%% ===================================================================
\partiel{Lexique à mémoriser}
\begin{multicols}{2}\small
\begin{itemize}[label=--,itemsep=0.8mm,leftmargin=*]
  \item \textbf{a doorway} — une embrasure de porte
  \item \textbf{flat} — plat
  \item \textbf{the edge} — le bord
  \item \textbf{dull} — terne
  \item \textbf{a cheek} — une joue
  \item \textbf{thin} — mince, maigre
  \item \textbf{to smile} — sourire
  \item \textbf{laughter} — le rire
  \item \textbf{merry} — joyeux
  \item \textbf{a beard} — une barbe
  \item \textbf{rough} — rugueux, grossier
  \item \textbf{surroundings} — l'entourage, le décor
\end{itemize}
\end{multicols}

\end{document}
